% --- Avant-propos -----------------------------------------------------------
\chapter*{Avant-propos}
\markboth{Avant-propos}{Avant-propos}
\addcontentsline{toc}{chapter}{Avant-propos}

Ce livre est né d'un constat simple : entre le programme officiel, qui dit
\emph{ce qu'il faut} enseigner, et la salle de classe, où il faut décider
\emph{comment} le faire tenir en une heure devant quarante élèves, il manque
un document. Les répartitions annuelles donnent le calendrier, les documents
d'accompagnement donnent l'esprit, mais aucun des deux ne met entre les mains
du professeur et de l'élève le même objet, ouvert à la même page.

C'est ce qu'on a essayé de faire ici. Chaque chapitre est écrit deux fois dans
un seul texte : l'élève y trouve un cours qu'il peut lire seul le soir, le
professeur y trouve les corrigés, les erreurs qu'il verra passer et les
remarques de conduite de séance. Les deux éditions sortent du même fichier
source ; rien ne peut donc diverger entre ce que l'un explique et ce que
l'autre révise.

\section*{Un mot à l'élève de terminale}

Vous entrez dans l'année où les mathématiques cessent d'être une succession
de techniques pour devenir un raisonnement continu. La dérivée de septembre
sert à l'intégrale de novembre, qui sert à l'équation différentielle de
janvier. Rien ne se referme derrière vous.

Vous remarquerez que chaque chapitre commence par une page de révision qui
ne parle jamais de la notion nouvelle. C'est voulu. On ne vous demandera
jamais de deviner ce que vous n'avez pas encore appris : on vérifie d'abord
que les outils de l'an dernier répondent encore, puis on ouvre le cours.
Prenez ces pages au sérieux, elles coûtent dix minutes et vous en font gagner
beaucoup.

\section*{Un mot au professeur}

Les volumes horaires indiqués en tête de chaque chapitre sont ceux du
Programme d'Études : soixante heures pour l'analyse, quarante-cinq pour
chacune des trois autres composantes. Ils sont serrés. Le tableau de
progression annuelle, quelques pages plus loin, propose une répartition
compatible avec les cinq périodes du calendrier scolaire, et indique pour
chaque chapitre ce qui peut être allégé si la période est amputée par les
examens blancs ou les intempéries.

Les démonstrations sont toutes écrites. Cela ne veut pas dire qu'elles doivent
toutes être faites au tableau : à vous de choisir. Elles sont là pour que
l'élève qui veut comprendre pourquoi un théorème est vrai trouve la réponse
dans son livre plutôt que nulle part.

\vspace{1.2em}
\noindent\textit{Antananarivo}
